\section{MLP-Mixer：进一步削弱归纳偏置的对照实验}

讲座最后用 MLP-Mixer 延伸“让数据说话”的思想。Mixer 不使用 convolution，也不使用 self-attention，而是在 token 维与 channel 维交替应用 MLP。它不是要证明 attention 无用，而是测试：在大规模预训练下，简单 dense mixing 能否学到足够好的视觉表示。

\subsection{从 ViT 到 Mixer：研究问题而非品牌替换}

过渡页把 ViT 的成功重新解释为“也许规模与数据比特定 operator 更重要”。因此 Mixer 是一个 controlled provocation：保留 patch embedding、residual/normalization 与大规模 recipe，替换 token interaction mechanism，观察 transfer scaling 是否仍成立。

\lecturefigure{slide-34-mixer-divider.jpg}{MLP-Mixer 延续“letting the data speak”，进一步移除 attention。}{官方视频 1:03:36。}

\begin{knowledgebox}{读图：不要得出“架构不重要”}
若 Mixer 在大数据上有竞争力，只能说明 attention 不是获得强表示的唯一机制。训练稳定性、compute efficiency、resolution scaling、dense prediction 与小数据性能仍可能不同。更合理的结论是：operator、scale 与 recipe 存在多种可行组合，需要按任务和系统约束选择。
\end{knowledgebox}

\subsection{Token-mixing 与 channel-mixing}

输入矩阵 $X\in\mathbb{R}^{N\times D}$ 中，$N$ 是 patch tokens，$D$ 是 channels。Token-mixing MLP 对每个 channel 跨 $N$ 个位置混合；channel-mixing MLP 对每个 token 跨 $D$ 个通道混合。两者交替并配合 residual connection，使信息既能跨空间位置传播，也能在特征维组合。

\lecturefigure{slide-35-mlp-mixer-architecture.jpg}{MLP-Mixer 交替执行 token-mixing 与 channel-mixing MLP。}{官方视频 1:03:44；Tolstikhin et al. 2021。}

\begin{importantbox}{Mixer block 的抽象形式}
忽略 LayerNorm 细节，可写为
\[
U=X+W_2\,\sigma(W_1X^{\top})^{\top},\qquad
Y=U+\sigma(UV_1)V_2.
\]
第一式沿 token 维混合，第二式沿 channel 维混合；$\sigma$ 是非线性。与 attention 不同，token mixing 权重不依赖当前输入内容，因此交互是静态的；代价和泛化特性也随 token 数与权重矩阵大小变化。
\end{importantbox}

\subsection{Scaling comparison：模型族的斜率比单点更重要}

最后一张图比较 Mixer、ViT 与 BiT/ResNet 随训练数据增长的 linear few-shot 表现。重点是三类模型都能从规模获益，但斜率、低数据起点与高数据上限不同；这再次支持“架构先验决定 sample efficiency，数据规模决定灵活模型能否追上”的框架。

\lecturefigure{slide-36-mixer-scaling.jpg}{Mixer、ViT 与 BiT 在不同训练集规模下的 linear few-shot scaling。}{官方视频 1:05:32。}

\begin{knowledgebox}{读图：比较起点、斜率和饱和区}
低数据端，强视觉先验模型通常有更高起点；随着训练数据增加，ViT/Mixer 曲线可能更陡；高数据端是否继续领先取决于容量与 recipe。单个最大数据点不能说明 sample efficiency，单个最小数据点也不能说明 asymptotic capability。完整曲线才显示 inductive bias 与 scale 的交换关系。
\end{knowledgebox}

\teachervoice{讲者把 Mixer 放在结尾，不是宣布“MLP 取代 Transformer”，而是强调研究应继续拆解：究竟是 attention、patch representation、规模、数据还是训练制度贡献了结果。老师的口头态度比 slide 标题更谨慎。}

\subsection{本章小结}

MLP-Mixer 证明 attention 不是大规模视觉表示的唯一候选。它把讲座主线推到更一般的结论：归纳偏置、数据规模和系统效率之间存在交换，模型族应通过完整 scaling curve 与多任务 transfer 比较。

